%=============================================================================
% EXPLICIT DERIVATION CHAIN: alpha^{-1} = 137.03599985 FROM DFD
% Extracted from the DFD Unified Theory (v108) and Ab Initio v12 (k_max=60)
% Updated 2026-03-22 with cross-reference to Ab Initio v12 paper
%=============================================================================
\documentclass[12pt,a4paper]{article}
\usepackage{amsmath,amssymb,amsthm}
\usepackage{booktabs}
\usepackage{geometry}
\geometry{margin=2.5cm}
\usepackage{hyperref}

\newtheorem{theorem}{Theorem}
\newtheorem{lemma}{Lemma}
\newtheorem{definition}{Definition}

\title{The Explicit DFD Derivation Chain for $\alpha^{-1} = 137.036$\\[6pt]
\large From Table~XXVIII Components to Sub-ppm Precision}
\author{Extracted from G.\ Alcock, \textit{Density Field Dynamics: A Complete Unified Theory}}
\date{March 2026}

\begin{document}
\maketitle

\begin{abstract}
We extract and present the complete derivation chain by which Density Field
Dynamics (DFD) obtains $\alpha^{-1} = 137.03599985$ (residual $-0.006$~ppm
from experiment) from eight locked geometric and Standard Model inputs.
The derivation has \textbf{two layers}: (1)~an exact electroweak mixing
formula $\alpha^{-1} = 6\pi\kappa_{U(1)}$ relating the fine-structure
constant to the $U(1)$ stiffness coefficient, and (2)~a spectral-action
computation on the Toeplitz-truncated $\mathbb{C}P^2 \times S^3$ internal
geometry that predicts $\kappa_{U(1)}$ from first principles.
We verify numerically that the lattice route reproduces $\alpha = 1/137$
within $\sim 1\%$, and explain why a circulating Reddit formula
``$\beta \times 6 \times (63/64) \times (4096/4095) \to 282$'' is incorrect.
\end{abstract}

\tableofcontents
\newpage

%=============================================================================
\section{The Eight Components (Table~XXVIII)}
\label{sec:components}
%=============================================================================

The convention-locked derivation of $\alpha^{-1}$ (Section~8C of the Unified
Theory) uses exactly eight inputs, all fixed by geometry or Standard Model
content:

\begin{table}[h]
\centering
\caption{Complete derivation chain for $\alpha^{-1}$
(Table~XXVIII of the Unified Theory).}
\label{tab:chain}
\renewcommand{\arraystretch}{1.3}
\begin{tabular}{llll}
\toprule
\textbf{Component} & \textbf{Value} & \textbf{Source} & \textbf{Status} \\
\midrule
$K_{\mathbb{C}P^2} = \mathcal{O}(-3)$ & $-3$
  & Algebraic geometry & Rigorous \\
$L_{\det} = K^{-1}$ & $\mathcal{O}(3)$
  & Spin$^c$ structure & Rigorous \\
$d = k_{\max} + 4$ & $64$
  & $\dim H^0(\mathcal{O}(k_{\max}+3))$ & Rigorous \\
$(d-1)/d$ & $63/64$
  & Traceless projection & Derived \\
$N_{\mathrm{species}}$ & $7$
  & SM SU(2) doublet components & SM content \\
$\mathrm{Tr}(Y^2)$ & $10$
  & SM hypercharge sum & SM content \\
$g_F$ & $8$
  & Spectral triple ($J \times \gamma \times \mathbb{C}$) & Derived \\
$w = N_{\mathrm{species}}/(g_F \cdot \mathrm{Tr}(Y^2))$ & $7/80$
  & Hypercharge weighting & Derived \\
$\varepsilon_{\mathrm{adj}} = d^2/(d^2-1)$ & $4096/4095$
  & Regular-module BOOST & Forced \\
$\beta_{U(1)}$ & $3.7969$
  & CS weighted average at $k_{\max}=60$ & Computed \\
\midrule
$\alpha^{-1}$ & $137.03599985$
  & All above combined & $< 0.01$~ppm \\
\bottomrule
\end{tabular}
\end{table}

%=============================================================================
\section{Layer 1: The Exact Electroweak Formula}
\label{sec:layer1}
%=============================================================================

The electromagnetic coupling emerges from mixing of $U(1)_Y$ and neutral
$SU(2)_L$ components. In the DFD gauge-emergence framework, the stiffness
functional for gauge sector $r$ is (Eq.~7 of Ref.~[1]):
\begin{equation}
\label{eq:stiffness}
\mathcal{L}^{(r)}_{\mathrm{stiff}}
  = -\frac{\kappa_r}{2}\,\mathrm{Tr}\,F^{(r)}_{ij} F^{(r)}_{ij}\,,
\end{equation}
where $\kappa_r = n_r \kappa_0$ with $n_{U(1)} = 1$, $n_{SU(2)} = 2$
(Frame Stiffness Theorem).  The gauge couplings are extracted via the
Wilson normalization:
\begin{equation}
\label{eq:wilson-norm}
g_1^2 = \frac{1}{\kappa_{U(1)}}\,,\qquad
g_2^2 = \frac{4}{\kappa_{SU(2)}}\,.
\end{equation}

The factor 4 in $g_2^2$ is the standard $SU(2)$ Wilson action normalization:
$\beta_{SU(2)} = 4/g_2^2$ from $\mathrm{Tr}(T^a T^b) = \tfrac{1}{2}\delta^{ab}$.

\paragraph{Electroweak mixing.}
The electromagnetic coupling $e$ satisfies (Eq.~10 of Ref.~[1]):
\begin{equation}
\label{eq:ew-mixing}
\frac{1}{e^2} = \frac{1}{g_1^2} + \frac{1}{g_2^2}
= \kappa_{U(1)} + \frac{\kappa_{SU(2)}}{4}\,.
\end{equation}

\paragraph{Applying the stiffness ratio.}
The Frame Stiffness Theorem gives $\kappa_{U(1)}/\kappa_{SU(2)} = n_1/n_2 = 1/2$,
so $\kappa_{SU(2)} = 2\kappa_{U(1)}$.  Substituting:
\begin{equation}
\frac{1}{e^2} = \kappa_{U(1)} + \frac{2\kappa_{U(1)}}{4}
= \frac{3}{2}\,\kappa_{U(1)}\,.
\end{equation}

The fine-structure constant is $\alpha = e^2/(4\pi)$, giving:
\begin{equation}
\boxed{
\alpha^{-1} = \frac{4\pi}{e^2}
= 4\pi \cdot \frac{3}{2}\,\kappa_{U(1)}
= 6\pi\,\kappa_{U(1)}\,.
}
\label{eq:master}
\end{equation}

This is the \textbf{exact electroweak formula}: the fine-structure constant
is determined entirely by the $U(1)$ stiffness coefficient.
Numerically, $\alpha^{-1}_{\mathrm{exp}} = 137.035999084$ requires:
\begin{equation}
\kappa_{U(1)} = \frac{137.035999084}{6\pi} = 7.269986\ldots
\end{equation}

%=============================================================================
\section{Layer 2: The Chern--Simons / Spectral Action Route}
\label{sec:layer2}
%=============================================================================

The question then reduces to: \emph{what predicts $\kappa_{U(1)}$?}
DFD provides two routes of different precision.

%---------------------------------------------------------------------
\subsection{Route A: Lattice Monte Carlo ($\sim 1\%$ precision)}
\label{sec:lattice}
%---------------------------------------------------------------------

The DFD microsector Chern--Simons level sum fixes the bare lattice coupling:
\begin{equation}
\label{eq:beta-def}
\beta_{U(1)} = \langle k_{\mathrm{eff}} \rangle
= \frac{\displaystyle\sum_{k=0}^{k_{\max}-1}(k+2)\,w(k)}
       {\displaystyle\sum_{k=0}^{k_{\max}-1} w(k)}\,,
\qquad
w(k) = \frac{2}{k+2}\sin^2\!\frac{\pi}{k+2}\,,
\end{equation}
where $k_{\max} = 60$ from the Bridge Lemma
($k_{\max} = \chi(\mathbb{C}P^2,\,\mathcal{O}(9)\oplus\mathcal{O}^{\oplus 5})
= 55 + 5 = 60$).

\paragraph{Numerical evaluation.}
\begin{equation}
\beta_{U(1)}\big|_{k_{\max}=60} = 3.7969\ldots \approx 3.80\,.
\end{equation}

The Wilson ratio $\beta_{SU(2)}/\beta_{U(1)} = (n_2/n_1) \times N_{\mathrm{gen}}
= 2 \times 3 = 6$ then gives $\beta_{SU(2)} = 22.80$.

\paragraph{Lattice measurement.}
At the derived parameter point $(\beta_{U(1)}, \beta_{SU(2)}) = (3.80, 22.80)$,
lattice Monte Carlo simulations \emph{measure} the renormalized stiffnesses
$\kappa_{U(1)}$ and $\kappa_{SU(2)}$ via the background-field method.
The fine-structure constant is then extracted using Eq.~\eqref{eq:master}.

\textbf{Key result} (86 lattice runs, $L = 4$--$16$):
\begin{equation}
\alpha_W = 0.007297 \pm 0.000094
\quad\Longrightarrow\quad
\alpha^{-1} = 137.0 \pm 1.8
\quad(\sim 1\%\text{ precision}).
\end{equation}

The crucial point: $\alpha$ was \emph{never used as input}.
The entire chain runs SM~$\to$~topology~$\to$~$\alpha$ (Eq.~A6 of Ref.~[1]).

%---------------------------------------------------------------------
\subsection{Route B: Spectral Action on Toeplitz-Truncated Geometry
(sub-ppm precision)}
\label{sec:spectral}
%---------------------------------------------------------------------

The full DFD derivation (Unified Theory, Section~8C) bypasses the lattice
and computes $\kappa_{U(1)}$ analytically from the $a_4$ Seeley--DeWitt
coefficient of the spectral action on $\mathbb{C}P^2 \times S^3$ with
Toeplitz truncation at level $d = k_{\max} + 4 = 64$.

The computation incorporates all eight Table~XXVIII components:

\begin{enumerate}
\item \textbf{Operator:} Connection Laplacian $P = -g^{ij}\nabla_i\nabla_j$
  on the internal bundle (K\"ahler form parallel $\Rightarrow$ no derivative
  ambiguities in higher $a_n$).

\item \textbf{Regulator:} Plateau cutoff with $f_{-2} = f_{-4} = \cdots = 0$
  (eliminates $a_6$ and higher as hidden tuning knobs).

\item \textbf{Finite-$k$ truncation:} Toeplitz quantization at level $m = k+3$
  on $\mathbb{C}P^1 \subset \mathbb{C}P^2$, giving $d = k+4 = 64$ and
  spectral cutoff $\Lambda^3 = k\,(d{-}1)/d = 60 \times 63/64 = 59.0625$.

\item \textbf{SM content:} The hypercharge weighting
  $w = N_{\mathrm{species}}/(g_F \cdot \mathrm{Tr}(Y^2)) = 7/80$
  enters through the $a_4$ trace over the internal fermion Hilbert space.

\item \textbf{Microsector trace:} The regular-module completion
  ($\mathcal{H}_F = M_d(\mathbb{C})$) forces the BOOST factor
  $\varepsilon_{\mathrm{adj}} = d^2/(d^2{-}1) = 4096/4095$
  from the democratic-to-physics trace conversion.
\end{enumerate}

\paragraph{The forced binary fork.}
Two microsector trace choices exist:
\begin{center}
\renewcommand{\arraystretch}{1.3}
\begin{tabular}{lccc}
\toprule
Branch & Factor & $\alpha^{-1}$ & Residual (ppm) \\
\midrule
\textbf{A (regular-module)} & $4096/4095$ & $137.03599985$ & $-0.006$ \\
B (fermion-rep) & $4095/4096$ & $137.03014445$ & $+42.7$ \\
\midrule
Experimental & --- & $137.035999084$ & --- \\
\bottomrule
\end{tabular}
\end{center}

Under the no-hidden-knobs policy, Branch~A (BOOST) survives; Branch~B
is falsified (43~ppm deficit, unrescuable).

\paragraph{Additive structure of the boost.}
The two branches decompose as:
\begin{equation}
\alpha^{-1} = F + G \cdot \delta\varepsilon\,,
\end{equation}
where $F \approx 137.033$, $G \approx 12$, and
$\delta\varepsilon = +1/(d^2{-}1)$ for Branch~A or
$\delta\varepsilon = -1/d^2$ for Branch~B.
The coefficient $G \approx 12 = 4 \times N_{\mathrm{gen}}$ reflects the
three-generation structure of the microsector trace.

%=============================================================================
\section{Why No Simple One-Liner Exists}
\label{sec:no-oneliner}
%=============================================================================

A key finding of this analysis: \textbf{there is no single closed-form
algebraic expression} $\alpha^{-1} = f(\beta, d, w, \varepsilon_{\mathrm{adj}})$
written in the DFD papers.  The reason is structural:

\begin{itemize}
\item The \textbf{lattice route} (Route~A) requires Monte Carlo simulation to
  measure the renormalized stiffness $\kappa_{U(1)}$ from the bare coupling
  $\beta_{U(1)}$.  The relationship $\kappa = \kappa(\beta)$ is non-perturbative
  and lattice-size-dependent.

\item The \textbf{spectral action route} (Route~B) computes $\kappa_{U(1)}$
  from the heat-kernel expansion of the connection Laplacian on the
  Toeplitz-truncated internal space.  This involves a numerical evaluation
  of the $a_4$ coefficient with all eight inputs, not a single algebraic
  ratio.
\end{itemize}

The closest to a ``one-liner'' is the \emph{exact} electroweak formula
\begin{equation}
\boxed{\alpha^{-1} = 6\pi\,\kappa_{U(1)}}
\end{equation}
combined with the \emph{numerical} prediction $\kappa_{U(1)} = 7.26999\ldots$
from the spectral action.

\paragraph{Approximate closed form.}
The ratio $\alpha^{-1}/(4\pi\,\beta\,(d{-}1)/d\cdot\varepsilon_{\mathrm{adj}})
\approx 2.9169$ is numerically close to $35/12 \approx 2.9167$
(within 85~ppm), suggesting an approximate formula:
\begin{equation}
\alpha^{-1} \approx \frac{35\pi}{3}\;\beta_{U(1)}\;\frac{d-1}{d}
\;\frac{d^2}{d^2-1}\,,
\qquad\text{(approximate, 85~ppm residual)}
\end{equation}
This factors as:
\begin{equation}
\frac{35\pi}{3} = 4\pi \times \underbrace{\frac{5}{2}}_{\text{EW mixing}} \times \underbrace{\frac{7}{6}}_{\text{spectral triple}}
\end{equation}
where $5/2 = 1 + R_W/4$ with Wilson ratio $R_W=6$ encodes the electroweak
$1/e^2 = 1/g'^2 + 1/g^2 = \beta + 3\beta/2 = 5\beta/2$, and
$7/6 = N_{\mathrm{species}}/(N_{\mathrm{gen}} \cdot n_2) = 7/(3 \times 2)$
encodes the microsector spectral triple correction.
This is \emph{not} the exact sub-ppm formula; the full spectral action
computation is required for that precision.

%=============================================================================
\section{Why ``$\beta \times 6 \times (63/64) \times (4096/4095)$'' Is Wrong}
\label{sec:wrong}
%=============================================================================

\paragraph{Numerical evaluation.}
\begin{equation}
\beta_{U(1)} \times 6 \times \frac{63}{64} \times \frac{4096}{4095}
= 3.7969 \times 6 \times 0.984375 \times 1.000244
= 22.43\,.
\end{equation}

This equals $\beta_{SU(2)} \times (63/64) \times (4096/4095)$---the
SU(2) bare lattice coupling with tiny traceless and boost corrections---\emph{not}
$\alpha^{-1}$.  The result 22.43 is neither 137 nor any other physically
meaningful quantity.

\paragraph{Four independent errors.}
\begin{enumerate}
\item \textbf{Missing $4\pi$ factor.}
  The fine-structure constant is defined as $\alpha = e^2/(4\pi)$, so
  $\alpha^{-1}$ must contain a factor of $4\pi \approx 12.566$.
  The naive expression omits this entirely.

\item \textbf{Misuse of the Wilson ratio.}
  The Wilson ratio $R_W = \beta_{SU(2)}/\beta_{U(1)} = 6$ determines the
  \emph{ratio} of lattice couplings.  It enters $\alpha^{-1}$ through
  electroweak mixing via
  \begin{equation}
  \frac{1}{e^2} = \frac{1}{g'^2} + \frac{1}{g^2}
    = \beta_{U(1)} + \frac{\beta_{SU(2)}}{4}
    = \beta_{U(1)}\!\left(1 + \frac{R_W}{4}\right)
    = \frac{5}{2}\,\beta_{U(1)},
  \end{equation}
  producing a factor of $5/2$, \textbf{not} a factor of $6$.
  Even the tree-level formula $\alpha^{-1}_{\mathrm{tree}} = 4\pi \times
  (5/2) \times \beta_{U(1)} = 10\pi\beta \approx 119.3$ differs radically
  from $6\beta \approx 22.4$.

\item \textbf{Missing spectral triple correction.}
  The hypercharge weight $w = 7/80$, spectral grading $g_F = 8$, and
  hypercharge trace $\mathrm{Tr}(Y^2)=10$ encode the finite Hilbert space
  structure.  These enter through the spectral action $a_4$ coefficient and
  produce an additional correction factor $\mathcal{C}_{\mathrm{spectral}}
  \approx N_{\mathrm{species}}/(N_{\mathrm{gen}} \cdot n_2) = 7/6$ that the
  naive formula omits.

\item \textbf{The result is $\beta_{SU(2)}$, not $\alpha^{-1}$.}
  The naive expression computes $6 \times \beta_{U(1)} \times
  (\text{tiny corrections}) \approx \beta_{SU(2)} = 22.8$, which is merely
  a bare SU(2) lattice coupling---an \emph{input parameter}, not the
  electromagnetic fine-structure constant.
\end{enumerate}

\paragraph{Comparison of formulas.}
\begin{center}
\renewcommand{\arraystretch}{1.4}
\begin{tabular}{lrl}
\toprule
\textbf{Formula} & \textbf{Value} & \textbf{Status} \\
\midrule
$\beta \times 6 \times (63/64) \times (4096/4095)$ & $22.43$ & \textbf{Wrong} (bare $\beta_{SU(2)}$) \\
$10\pi\,\beta$ & $119.3$ & Tree-level; 13\% low \\
$(35\pi/3)\,\beta\,(63/64)\,(4096/4095)$ & $137.024$ & Approx.; 85~ppm residual \\
$6\pi\,\kappa_{U(1)}$ (spectral action) & $137.03599985$ & \textbf{Correct} ($-0.006$~ppm) \\
\bottomrule
\end{tabular}
\end{center}

%=============================================================================
\section{Complete Numerical Verification}
\label{sec:numerics}
%=============================================================================

%---------------------------------------------------------------------
\subsection{Computing $\beta_{U(1)}$}
%---------------------------------------------------------------------

With $k_{\max} = 60$:
\begin{equation}
\beta_{U(1)} = \frac{\sum_{k=0}^{59}(k+2)\,w(k)}{\sum_{k=0}^{59}w(k)}
= \frac{8.3245}{2.1924} = 3.796945\ldots
\end{equation}

\noindent Python verification:
\begin{verbatim}
import math
def w(k): return (2.0/(k+2)) * math.sin(math.pi/(k+2))**2
Z = sum(w(k) for k in range(60))
k_eff = sum((k+2)*w(k) for k in range(60)) / Z
# k_eff = 3.796945275...
\end{verbatim}

%---------------------------------------------------------------------
\subsection{Lattice-route verification ($\sim 1\%$)}
%---------------------------------------------------------------------

From L12 lattice data: $\kappa_{U(1)} \approx 7.27$, $\kappa_{SU(2)} \approx 14.7$.

Using Eq.~\eqref{eq:master}:
\begin{equation}
\alpha^{-1} = 6\pi \times 7.27 = 137.1\,.
\end{equation}

More precisely, the L12 result $\alpha_W = 0.007291$ gives
$\alpha^{-1} = 137.15$ ($-0.08\%$ from experiment).

%---------------------------------------------------------------------
\subsection{Spectral-action verification (sub-ppm)}
%---------------------------------------------------------------------

The spectral action on the Toeplitz-truncated $\mathbb{C}P^2 \times S^3$
with the eight locked components gives:

\begin{equation}
\kappa_{U(1)}^{\mathrm{pred}} = 7.26999\ldots
\end{equation}

Therefore:
\begin{equation}
\alpha^{-1} = 6\pi \times 7.26999\ldots = 137.03599985\,.
\end{equation}

\noindent Residual from experiment:
\begin{equation}
\frac{\alpha^{-1}_{\mathrm{DFD}} - \alpha^{-1}_{\mathrm{exp}}}
     {\alpha^{-1}_{\mathrm{exp}}}
= \frac{137.03599985 - 137.035999084}{137.035999084}
= -0.006\text{ ppm}\,.
\end{equation}

%=============================================================================
\section{The Derivation Chain in One Diagram}
\label{sec:diagram}
%=============================================================================

\begin{center}
\fbox{\parbox{0.9\textwidth}{
\textbf{SM hypercharge spectrum} $\to$ $q_1 = 3$ (anomaly cancellation)

$\downarrow$

$a = q_1^2 \cdot q_1 = 9$ $\to$ twist bundle $E = \mathcal{O}(9) \oplus \mathcal{O}^{\oplus 5}$

$\downarrow$

$k_{\max} = \chi(\mathbb{C}P^2, E) = 55 + 5 = 60$ \quad [Bridge Lemma]

$\downarrow$

$\beta_{U(1)} = \langle k+2 \rangle_{k_{\max}=60} = 3.7969$ \quad [CS weighted average]

$\downarrow$ \quad + \quad Wilson ratio $= 6$ \quad + \quad $\kappa_{U(1)}/\kappa_{SU(2)} = 1/2$

$\downarrow$

\textbf{Lattice route:} MC at $(\beta_{U(1)}, \beta_{SU(2)}) = (3.80, 22.80)$

$\to$ measure $\kappa_{U(1)} \approx 7.27$ $\to$ $\alpha^{-1} = 6\pi\kappa \approx 137$ ($\pm 1\%$)

$\downarrow$ \quad OR

\textbf{Spectral action route:} $a_4$ on Toeplitz-truncated $\mathbb{C}P^2 \times S^3$

+ traceless $(63/64)$ + BOOST $(4096/4095)$ + SM content $(N_{\mathrm{species}}, g_F, \mathrm{Tr}(Y^2))$

$\to$ $\kappa_{U(1)} = 7.26999\ldots$ $\to$ $\alpha^{-1} = 137.03599985$ (sub-ppm)
}}
\end{center}

%=============================================================================
\section{Summary}
\label{sec:summary}
%=============================================================================

\begin{enumerate}
\item The \textbf{exact electroweak formula} is:
\begin{equation}
\alpha^{-1} = 6\pi\,\kappa_{U(1)}\,,
\end{equation}
derived from $1/e^2 = 1/g_1^2 + 1/g_2^2$ with the DFD stiffness ratio
$\kappa_{U(1)}/\kappa_{SU(2)} = 1/2$.

\item The stiffness $\kappa_{U(1)}$ is \textbf{predicted} from the eight
  Table~XXVIII components via the spectral action on the Toeplitz-truncated
  internal geometry.

\item The sub-ppm result $\alpha^{-1} = 137.03599985$ comes from a
  \textbf{numerical evaluation} of the $a_4$ Seeley--DeWitt coefficient,
  not from a single closed-form algebraic expression.

\item The \textbf{lattice Monte Carlo} independently confirms $\alpha = 1/137$
  within $\sim 1\%$ at the derived parameter point.

\item The Reddit formula ``$\beta \times 6 \times (63/64) \times (4096/4095)$''
  gives $22.43$, not $137$ or $282$.  It confuses the Wilson ratio (which
  relates lattice couplings) with a direct multiplier in the $\alpha$ formula,
  and omits the electroweak mixing, stiffness renormalization, and SM trace
  factors.

\item All inputs are topologically or SM-fixed; there are
  \textbf{zero continuous free parameters}.
\end{enumerate}

%=============================================================================
\section*{References}
%=============================================================================

\begin{enumerate}
\item[{[1]}] G.~Alcock, ``Ab Initio Derivation of the Fine Structure Constant
  from Density Field Dynamics,'' (2025).
  \url{https://doi.org/10.5281/zenodo.19173548}

\item[{[2]}] G.~Alcock, ``Density Field Dynamics: A Complete Unified Theory,''
  v108, Zenodo (2026).  Section~8C (convention-locked $\alpha$), Appendix~K
  (microsector physics), Appendix~F (gauge emergence).
\end{enumerate}

\end{document}
